\documentclass[10pt,a4paper]{article}
\usepackage{cvcommon}
\hypersetup{pdftitle={CV Javier López Molinero - Español},pdfauthor={Javier López Molinero},pdfsubject={Currículum profesional}}

\begin{document}
\CVHeader{\HeaderContent{SENIOR C++ / QUANT DEVELOPER\enspace·\enspace DEVOPS}{Madrid, España}{Más de 15 años en software, integración y automatización}}
\RightPanel{6.08cm}
\columnratio{0.686}\setlength{\columnsep}{1.06cm}
\begin{paracol}{2}
\vspace{-0.20cm}
\CVSection{Perfil}
{\fontsize{8.25}{10.8}\selectfont
Ingeniero de Telecomunicaciones con más de 15 años de experiencia en desarrollo de software, integración de sistemas e infraestructura CI/CD. Especializado en C++, Python y Java, con enfoque DevOps, sistemas de tiempo real, microservicios y entornos Linux.}\par
\vspace{0.30cm}

\CVSection{Experiencia profesional}
\FeaturedJob{Quant Developer SR and Integrator Engineer for BBVA}{ARENA Financial Tech}{Abr. 2026 - Actualidad}{Desarrollo de software y mejora de la infraestructura de integración continua de BBVA en la Sala de Tesorería.}
% METRICA: añadir reducción del tiempo de compilación, frecuencia de despliegue o incidencias evitadas.
\Job{Ingeniero desarrollador e integrador para BBVA}{Optimissa-Alten}{Mar. 2022 - Mar. 2026}{Desarrollo en C++ y Python y tareas DevOps para Quants Develop Design and Cloud en la Sala de Tesorería.}{}
\Job{Ingeniero desarrollador e integrador para Telefónica}{Full On Net}{Nov. 2019 - Mar. 2022}{CI/CD para IPTV Movistar: C++, Python, servidores, redes, migración de entornos y Yocto Linux.}{}
\Job{Advanced Development Engineer}{Indra · Proyecto HERMES / HORUS}{Jul. 2018 - Nov. 2019}{Centro de control urbano: C++, Java, Python, microservicios, Docker, Kubernetes, DevOps y scripting.}{}
\Job{Ingeniero de Desarrollo}{Altran · Cliente Ericsson}{Oct. 2016 - Jul. 2018}{Equipo cross-functional de I+D; desarrollo en C++ y TTCN-3 e integración continua para el nodo HSS.}{}
\Job{Ingeniero de Software}{Indra Sistemas / Novanotio}{Feb. 2012 - Oct. 2016}{Firmware para sistemas de control de tráfico vial y coordinación del desarrollo del software del radar Doppler.}{}

\switchcolumn
\vspace{0.42cm}
\SideSection{Tecnologías}{
  \item C · C++ · Python · Java · ADA · PHP
  \item Docker · Kubernetes · Jenkins · Gerrit
  \item CI/CD · microservicios · testing · Scrum
  \item Linux · Windows Server · macOS · Unix
}
\SideSection{Especialización}{
  \item Desarrollo de software con C++
  \item DevOps e integración de sistemas
  \item Cloud, SaaS e IoT
  \item Tiempo real y firmware embebido
  \item TCP/IP · OSI · SDN · NFV
}
\SideSection{Fortalezas}{
  \item Gestión de equipos de desarrollo
  \item Resolución de problemas complejos
  \item Preventa, soporte y formación técnica
}
\SideSection{Idiomas}{
  \item Inglés: competencia profesional en entornos internacionales
}
\end{paracol}
\CVFooter{1}{2}

\newpage
\PageTwoHeader{PROYECTOS, FORMACIÓN Y TRAYECTORIA}{Perfil técnico senior · software, integración, automatización e infraestructuras}
\RightPanel{3.49cm}
\columnratio{0.686}\setlength{\columnsep}{1.06cm}
\begin{paracol}{2}
\vspace{-0.20cm}
\CVSection{Formación}
\EducationBox{
  \EduRow{Máster en Ingeniería de Telecomunicaciones}{UOC · 2014-2018}
  \EduRow{Ingeniería Técnica de Telecomunicaciones - Telemática}{Universidad de Alcalá · 2003-2009}
  \EduRow{Técnico Superior en Sistemas de Telecomunicaciones e Informáticos}{La Salle-Sagrado Corazón · 2001-2003}
}

\CVSection{Proyectos seleccionados}
\Project{BBVA - Sala de Tesorería}{Desarrollo de módulos y mejora continua de la infraestructura CI/CD para el área de Quants.}
\Project{Movistar IPTV}{Creación del entorno CI/CD y de la jaula de compilación para descodificadores, con C++, Python y Yocto Linux.}
\Project{HERMES / HORUS}{Back-end del centro de control semafórico de Indra con C++, Java y servicios distribuidos.}
\Project{Nodo HSS - Ericsson I+D}{Microservicios, integración continua y Kubernetes dentro de un equipo cross-functional.}
\Project{Control de tráfico}{Firmware para radares y periféricos; coordinación del desarrollo del software del radar Doppler.}
\Project{CERT del Gobierno de Venezuela}{Infraestructura de seguridad sobre nube privada persistente con servicios de correo, web, OTRS y OSSIM.}
\Project{IDS - Palacio de la Moncloa}{Implantación de una maqueta y procedimiento reproducible para la red interna.}
\Project{Ministerio de Defensa}{Aula de seguridad informática basada en máquinas virtuales y herramientas de ataque y defensa.}
\Project{Routers Funkwerk / Teldat}{Soporte nivel 2 y homologación para Portugal Telecom, Selta y Telefónica.}
\Project{Red corporativa de Teydisa}{Reestructuración integral de cableado, red e intranet.}
\Project{Banco Santander}{Mantenimiento y desarrollo de aplicaciones para la plataforma Moodle de másteres online.}
\Project{Radio online laisla.fm}{Diseño, creación e implantación en el mercado de la plataforma de radio online, incluido un reproductor de streaming para iPhone con más de 1.200 descargas.}

\switchcolumn
\vspace{0.42cm}
\SideSection{Formación complementaria}{
  \item 2012 · Programación Java (250 h)
  \item 2011 · CCNA, ESFgroup
  \item 2007 · Cisco CCNA 1 y 2, University of NEWI (Wales)
}
\SideSection{Experiencia anterior}{
  \item \textbf{TB-Security}\\Consultor Junior en Sistemas · abr.-ago. 2011
  \item \textbf{CIFF - Banco Santander}\\Gestión tecnológica y soporte Moodle · nov. 2010-abr. 2011
  \item \textbf{Funkwerk EC Iberia}\\Preventa, soporte, comercial y formación · mar.-may. 2010
  \item \textbf{Quantum 2003, S.L.}\\Servicios IT y telecomunicaciones · jun. 2004-oct. 2009
}
\SideSection{Proyecto fin de carrera}{
  \item Diseño, creación e implantación en el mercado de la radio online laisla.fm, incluida su plataforma tecnológica y el reproductor para iPhone.
}
\SideSection{Proyecto fin de máster}{
  \item \textbf{MQTT-S sobre IEEE 802.15.4e en OpenMote}. UOC, 2018. Prueba de concepto IoT con OpenMote, Raspberry Pi y sensores. Transmisión de datos mediante MQTT-S a una base de datos en la nube, visualización con Grafana y ampliación de OpenWSN para evaluar la viabilidad del protocolo.
}
\SideSection{Aficiones}{
  \item Programación y proyectos personales
  \item Domótica y automatización del hogar
  \item Deporte, especialmente karate
  \item IA para mejorar procesos, aprender y explorar nuevas aplicaciones
}
\end{paracol}
\CVFooter{2}{2}
\end{document}
